% ============================================================
% Phase A - Réseau plat : préparation et attaques de référence
% Steps 1 à 4 du projet "Zero Trust Lab"
% À intégrer dans le modèle rapportENSIAS
% ============================================================

\chapter{Phase A --- Réseau plat : mise en place et attaques de référence}

\section{Introduction et objectifs}

Cette phase a pour objectif de construire une architecture volontairement
non sécurisée, représentative d'un déploiement rapide sans considération
particulière pour la sécurité réseau. Toutes les machines partagent le
même segment \texttt{192.168.1.0/24}, sans aucune règle de filtrage entre
elles. Cette configuration sert de \emph{baseline} : les vulnérabilités
qui y sont démontrées seront corrigées une à une lors de la Phase B
(architecture Zero Trust), permettant une comparaison avant/après
mesurable.

\section{Architecture réseau de la Phase A}

\begin{table}[h]
\centering
\begin{tabular}{|l|l|l|}
\hline
\textbf{Machine} & \textbf{Adresse IP} & \textbf{Rôle} \\
\hline
OPNsense & 192.168.1.1 & Routeur / pare-feu (sans filtrage actif) \\
Server-DB & 192.168.1.20 & Serveur de base de données MariaDB \\
Server-Nextcloud & 192.168.1.30 & Serveur applicatif Nextcloud (Apache/PHP) \\
Kali-Attacker & 192.168.1.x (DHCP) & Poste d'attaque \\
\hline
\end{tabular}
\caption{Adressage réseau de la Phase A (réseau plat)}
\label{tab:adressage-phase-a}
\end{table}

Toutes les machines sont reliées au même réseau interne VirtualBox
(\texttt{intnet-lan}), qui simule un commutateur (switch) Ethernet.
Aucune segmentation ni règle de pare-feu n'est appliquée à ce stade.

\section{Mise en place de Server-DB}

\subsection{Configuration réseau statique}

L'installateur Debian n'ayant configuré que l'interface utilisée pendant
l'installation (temporairement en NAT), l'interface réseau interne
\texttt{enp0s3} a dû être configurée manuellement en adresse statique :

\begin{lstlisting}[language=bash, caption={Configuration statique de l'interface réseau}]
auto enp0s3
iface enp0s3 inet static
    address 192.168.1.20/24
    gateway 192.168.1.1
    dns-nameservers 192.168.1.1
\end{lstlisting}

Après application (\texttt{systemctl restart networking}), la connectivité
a été validée par trois tests de ping (passerelle, DNS public, nom de
domaine), tous réussis avec 0\,\% de perte de paquets.

\begin{figure}[h]
\centering
% \includegraphics[width=0.8\textwidth]{images/server-db-ip-pings.png}
\caption{Validation de la configuration réseau de Server-DB : adresse IP et tests de connectivité}
\label{fig:server-db-network}
\end{figure}

\textbf{Point de dépannage noté} : lors de la première tentative de
configuration, un adaptateur NAT temporaire laissé actif par
l'installateur permettait à la machine de contourner le pare-feu OPNsense
pour accéder à internet. Ce comportement a été identifié et corrigé en
désactivant l'adaptateur NAT, forçant toute la sortie réseau à transiter
par OPNsense.

\subsection{Installation et sécurisation de MariaDB}

MariaDB (version 11.8.6) a été installé puis sécurisé à l'aide du script
\texttt{mariadb-secure-installation} :

\begin{itemize}
    \item suppression des utilisateurs anonymes ;
    \item interdiction de la connexion root à distance ;
    \item suppression de la base de données \texttt{test} ;
    \item rechargement des tables de privilèges.
\end{itemize}

\subsection{Exposition volontaire de la base de données}

Afin de reproduire une erreur de configuration courante (base de données
accessible depuis tout le réseau plutôt que restreinte au serveur
applicatif), deux modifications ont été appliquées :

\begin{lstlisting}[language=SQL, caption={Création de la base et de l'utilisateur Nextcloud}]
CREATE DATABASE nextcloud CHARACTER SET utf8mb4 COLLATE utf8mb4_general_ci;
CREATE USER 'nextcloud'@'192.168.1.%' IDENTIFIED BY 'DB_PASSWORD_REDACTED';
GRANT ALL PRIVILEGES ON nextcloud.* TO 'nextcloud'@'192.168.1.%';
FLUSH PRIVILEGES;
\end{lstlisting}

\begin{lstlisting}[language=bash, caption={Ouverture de MariaDB sur le réseau}]
# /etc/mysql/mariadb.conf.d/50-server.cnf
bind-address = 0.0.0.0
\end{lstlisting}

La vérification \texttt{ss -tlnp | grep 3306} confirme que le service
écoute sur \texttt{0.0.0.0:3306}, c'est-à-dire sur toutes les interfaces
réseau, et non uniquement en local. Cette exposition volontaire sera
exploitée au \S\,\ref{sec:attaque-db} et corrigée en Phase B par une
règle de pare-feu restreignant l'accès au seul serveur applicatif.

\begin{figure}[h]
\centering
% \includegraphics[width=0.8\textwidth]{images/server-db-bind-address.png}
\caption{MariaDB en écoute sur toutes les interfaces réseau (port 3306)}
\label{fig:mariadb-expose}
\end{figure}

\section{Mise en place de Server-Nextcloud}

\subsection{Configuration de la machine}

Une VM Debian 13 (2\,Go de RAM, 20\,Go de disque) a été créée avec un
unique adaptateur réseau en mode \emph{Internal Network}
(\texttt{intnet-lan}), sans adaptateur NAT, afin que l'installation des
paquets transite exclusivement par OPNsense. L'adresse statique
\texttt{192.168.1.30/24} a été appliquée selon la même méthode que pour
Server-DB.

\subsection{Pile applicative}

Apache 2.4.68 et PHP 8.4.26 ont été installés avec les extensions
requises par Nextcloud (\texttt{mysql}, \texttt{gd}, \texttt{curl},
\texttt{mbstring}, \texttt{xml}, etc.). La connectivité vers Server-DB a
été validée avant l'installation de Nextcloud :

\begin{lstlisting}[language=bash, caption={Test de connexion à la base de données distante}]
mariadb -h 192.168.1.20 -u nextcloud -p
\end{lstlisting}

La connexion réussie confirme qu'aucune restriction réseau n'empêche
Server-Nextcloud (ni, potentiellement, toute autre machine du segment)
d'atteindre la base de données.

\subsection{Installation de Nextcloud}

Nextcloud (dernière version stable) a été téléchargé, extrait dans
\texttt{/var/www/nextcloud}, puis installé en ligne de commande via
\texttt{occ}, sans interface web :

\begin{lstlisting}[language=bash, caption={Installation de Nextcloud via occ}]
sudo -u www-data php /var/www/nextcloud/occ maintenance:install \
    --database mysql \
    --database-name nextcloud \
    --database-user nextcloud \
    --database-pass DB_PASSWORD_REDACTED \
    --database-host 192.168.1.20 \
    --admin-user admin \
    --admin-pass ADMIN_PASSWORD_REDACTED
\end{lstlisting}

\begin{figure}[h]
\centering
% \includegraphics[width=0.8\textwidth]{images/nextcloud-dashboard.png}
\caption{Tableau de bord Nextcloud après installation réussie}
\label{fig:nextcloud-dashboard}
\end{figure}

\subsection{Comptes utilisateurs de test}

Afin de disposer de cibles réalistes pour les tests de brute force
(\S\,\ref{sec:bruteforce}), plusieurs comptes ont été créés avec des
mots de passe représentatifs de pratiques faibles courantes en
entreprise (nom dérivé, réutilisation de mots courants) :

\begin{table}[h]
\centering
\begin{tabular}{|l|l|l|}
\hline
\textbf{Utilisateur} & \textbf{Mot de passe} & \textbf{Pattern illustré} \\
\hline
admin & ADMIN_PASSWORD_REDACTED & Compte administrateur \\
admin2 & Qdmin123456789 & Faux compte admin secondaire oublié \\
jdupont & Jdupont123 & Mot de passe dérivé du nom \\
m.martin & ;qrtin;123456 & Mot de passe quasi-générique \\
y.alaoui & Alaoui12345678 & Mot de passe dérivé du nom \\
\hline
\end{tabular}
\caption{Comptes Nextcloud créés pour les tests d'intrusion}
\label{tab:comptes-nextcloud}
\end{table}

\textbf{Observation} : Nextcloud a automatiquement rejeté plusieurs mots
de passe initialement proposés (\texttt{Password1},
\texttt{Alaoui1234}) car présents dans une liste de mots de passe
compromis (vérification de type \emph{Have I Been Pwned}). Ce contrôle
natif constitue une première ligne de défense présente dès la Phase A,
à mentionner comme bonne pratique déjà en place.

\section{Step 4 --- Attaques de référence (baseline)}

Toutes les attaques suivantes ont été menées depuis une machine Kali
Linux dédiée (\texttt{Kali-Attacker}), clonée d'une image existante puis
reconfigurée sur le réseau \texttt{intnet-lan} uniquement, avec
régénération de son adresse MAC pour éviter tout conflit.

\subsection{Reconnaissance réseau (scan Nmap)}
\label{sec:nmap}

Un scan complet de tous les ports TCP, avec détection de version des
services, a été lancé sur l'ensemble du sous-réseau :

\begin{lstlisting}[language=bash, caption={Scan de reconnaissance Nmap}]
nmap -sV -p- 192.168.1.0/24
\end{lstlisting}

\begin{table}[h]
\centering
\begin{tabular}{|l|l|l|}
\hline
\textbf{Hôte} & \textbf{Ports ouverts} & \textbf{Services identifiés} \\
\hline
192.168.1.1 & 53, 80, 443, 53053 & OPNsense (Unbound DNS, interface web) \\
192.168.1.20 & 22, 3306 & OpenSSH 10.0, MariaDB (non filtrée) \\
192.168.1.30 & 22, 80 & OpenSSH 10.0, Apache 2.4.68 (HTTP) \\
\hline
\end{tabular}
\caption{Résultats du scan de reconnaissance}
\label{tab:nmap-resultats}
\end{table}

\textbf{Analyse} : le scan révèle, sans aucune authentification préalable,
l'intégralité des services exposés sur le réseau, y compris le port
\textbf{3306 (MySQL/MariaDB)} directement accessible depuis le poste
attaquant. Dans une architecture correctement segmentée, ce port ne
devrait être visible que depuis le serveur applicatif autorisé.

\begin{figure}[h]
\centering
% \includegraphics[width=0.9\textwidth]{images/nmap-scan-resultats.png}
\caption{Résultats du scan Nmap révélant l'ensemble des services exposés}
\label{fig:nmap-scan}
\end{figure}

\subsection{Chaîne d'attaque complète : de la reconnaissance à l'accès aux données}
\label{sec:attaque-db}

Plutôt que de supposer un accès direct et immédiat à la base de données
(ce qui supposerait que l'attaquant connaisse déjà les identifiants),
une chaîne d'attaque réaliste et complète a été reconstituée, respectant
la logique d'un attaquant ne disposant initialement d'aucune information
privilégiée.

\subsubsection{Étape 1 --- Collecte d'information (OSINT) et génération de dictionnaire ciblé}

En supposant que l'attaquant a pu identifier le nom de l'administratrice
du système (scénario réaliste : fuite d'information, réseau social,
métadonnées de document), un dictionnaire de mots de passe ciblé a été
généré avec l'outil \textbf{CUPP} (\emph{Common User Passwords
Profiler}), qui automatise la création de variantes probables à partir
d'informations personnelles (nom, prénom, suffixes numériques
courants) :

\begin{lstlisting}[language=bash, caption={Génération d'un dictionnaire ciblé avec CUPP}]
cupp -i
# Saisie interactive : prenom, nom, mots-cles
# Resultat : 11 940 variantes generees
\end{lstlisting}

\subsubsection{Étape 2 --- Attaque par force brute ciblée (SSH)}
\label{sec:bruteforce}

Un premier essai avec la liste complète (11\,940 mots) s'est révélé
impraticable en raison du mécanisme de limitation de connexions
simultanées de Debian (environ 22 à 24 heures estimées au rythme
observé). La liste a donc été réduite aux variantes les plus probables
(préfixes \texttt{imane} et \texttt{arrach}), ramenant le test à
quelques dizaines de mots :

\begin{lstlisting}[language=bash, caption={Attaque par force brute avec Hydra}]
hydra -l nextcloud-user -P imane_fast.txt -t 1 -W 2 ssh://192.168.1.30
\end{lstlisting}

\textbf{Résultat} : le mot de passe \texttt{imanearrach} a été trouvé en
10 tentatives, en moins de 2 minutes.

\begin{figure}[h]
\centering
% \includegraphics[width=0.9\textwidth]{images/hydra-ssh-success.png}
\caption{Succès de l'attaque par force brute ciblée contre le service SSH}
\label{fig:hydra-success}
\end{figure}

\subsubsection{Étape 3 --- Extraction des identifiants de base de données}

Une fois l'accès SSH obtenu, le fichier de configuration de Nextcloud a
été consulté. Comme pour la grande majorité des applications web, les
identifiants de connexion à la base de données y sont stockés en clair :

\begin{lstlisting}[language=bash, caption={Lecture du fichier de configuration}]
sudo cat /var/www/nextcloud/config/config.php
\end{lstlisting}

\begin{lstlisting}[language=PHP, caption={Extrait révélant le mot de passe en clair}]
'dbhost' => '192.168.1.20',
'dbuser' => 'nextcloud',
'dbpassword' => 'DB_PASSWORD_REDACTED',
\end{lstlisting}

\begin{figure}[h]
\centering
% \includegraphics[width=0.9\textwidth]{images/config-php-leak.png}
\caption{Mot de passe de base de données exposé en clair dans config.php}
\label{fig:config-leak}
\end{figure}

\subsubsection{Étape 4 --- Accès à la base de données et exfiltration}

Avec les identifiants extraits, la connexion directe à MariaDB (toujours
depuis Kali, démontrant la portée de l'exposition réseau identifiée au
\S\,\ref{sec:nmap}) a été établie sans aucune restriction :

\begin{lstlisting}[language=bash, caption={Connexion à la base de données avec les identifiants extraits}]
mariadb -h 192.168.1.20 -u nextcloud -p
USE nextcloud;
SHOW TABLES;
SELECT uid, password FROM oc_users;
\end{lstlisting}

La requête retourne l'intégralité des comptes Nextcloud et leurs
empreintes de mot de passe (hachage \textbf{Argon2id}, algorithme
résistant au cassage par calcul massivement parallèle, ce qui limite
l'impact immédiat de cette fuite mais démontre néanmoins une exposition
complète des données applicatives).

\begin{figure}[h]
\centering
% \includegraphics[width=0.9\textwidth]{images/mariadb-oc-users.png}
\caption{Accès complet à la table des utilisateurs Nextcloud depuis le poste attaquant}
\label{fig:mariadb-access}
\end{figure}

\subsubsection{Synthèse de la chaîne d'attaque}

\begin{table}[h]
\centering
\begin{tabular}{|c|l|l|}
\hline
\textbf{Étape} & \textbf{Technique} & \textbf{Résultat} \\
\hline
1 & Scan réseau (Nmap) & Découverte des services exposés \\
2 & OSINT + génération de dictionnaire (CUPP) & Liste de mots de passe ciblée \\
3 & Force brute SSH (Hydra) & Accès shell sur Server-Nextcloud \\
4 & Lecture de configuration & Mot de passe de base de données \\
5 & Connexion directe à la base & Exfiltration des comptes utilisateurs \\
\hline
\end{tabular}
\caption{Synthèse de la chaîne d'attaque complète}
\label{tab:synthese-attaque}
\end{table}

Cette chaîne illustre un principe central de la sécurité réseau :
\textbf{une seule porte ouverte (SSH avec mot de passe faible) suffit à
compromettre l'ensemble de la chaîne applicative}, dès lors qu'aucune
segmentation ni principe de moindre privilège n'est appliqué. C'est
précisément ce que la Phase B (architecture Zero Trust) se propose de
corriger.

\subsection{Interception du trafic réseau non chiffré}
\label{sec:sniffing}

\subsubsection{Limite du sniffing passif sur réseau commuté}

Une première tentative de capture du trafic entre \texttt{Client-Debian}
et \texttt{Server-Nextcloud} depuis Kali s'est révélée infructueuse (0
paquet capturé). Ce résultat, bien que négatif, constitue une
observation pédagogique importante : contrairement à un réseau basé sur
un concentrateur (\emph{hub}), un réseau commuté (\emph{switch}, tel que
simulé par le réseau interne VirtualBox) n'achemine le trafic qu'entre
l'émetteur et le destinataire concernés. Une machine tierce ne peut donc
intercepter passivement que son propre trafic, sauf à mettre en œuvre
une attaque active de type \emph{Man-in-the-Middle} (ARP spoofing), non
réalisée dans le cadre de ce test.

\subsubsection{Capture réussie}

La capture a été reproduite en initiant la connexion directement depuis
le poste attaquant :

\begin{lstlisting}[language=bash, caption={Capture du trafic HTTP}]
sudo tcpdump -i eth0 -w capture.pcap host 192.168.1.30 and port 80
\end{lstlisting}

Une connexion au compte \texttt{jdupont} a ensuite été effectuée depuis
le navigateur de Kali, pendant que la capture était active (1\,153
paquets capturés). L'analyse avec Wireshark (filtre
\texttt{http.request.method == "POST"}, puis \emph{Follow HTTP Stream})
révèle le contenu intégral de la requête de connexion :

\begin{lstlisting}[caption={Identifiants en clair interceptés dans le flux HTTP}]
POST /index.php/login HTTP/1.1
Host: 192.168.1.30
...
user=jdupont&password=Jdupont123&rememberme=1...
\end{lstlisting}

\begin{figure}[h]
\centering
% \includegraphics[width=0.9\textwidth]{images/wireshark-http-stream.png}
\caption{Nom d'utilisateur et mot de passe visibles en clair dans le flux HTTP intercepté}
\label{fig:wireshark-stream}
\end{figure}

\textbf{Analyse} : l'absence de chiffrement (HTTP au lieu de HTTPS) sur
la page d'authentification permet à tout observateur positionné sur le
chemin réseau d'obtenir les identifiants complets d'un utilisateur, sans
nécessiter de cassage de mot de passe. Cette vulnérabilité sera corrigée
en Phase B par la mise en place d'un reverse proxy Traefik imposant le
chiffrement TLS de bout en bout.

\section{Synthèse de la Phase A}

\begin{table}[h]
\centering
\begin{tabular}{|l|l|l|}
\hline
\textbf{Attaque} & \textbf{Statut} & \textbf{Vulnérabilité exploitée} \\
\hline
Scan réseau & Réussi & Absence de filtrage, tous les services visibles \\
Force brute SSH & Réussi (2 min) & Mot de passe faible et prévisible \\
Extraction de configuration & Réussi & Identifiants stockés en clair \\
Accès direct base de données & Réussi & Base de données non filtrée réseau \\
Mouvement latéral & Réussi & Absence de segmentation \\
Interception de trafic & Réussi (depuis l'hôte) & Absence de chiffrement TLS \\
\hline
\end{tabular}
\caption{Bilan des attaques de la Phase A --- baseline avant sécurisation}
\label{tab:bilan-phase-a}
\end{table}

Ce bilan constitue la ligne de référence \textit{« avant »} du tableau
de validation final (Step 12), qui sera complété après la mise en œuvre
de l'architecture Zero Trust en Phase B par la ligne \textit{« après »}
correspondante, permettant de mesurer objectivement la réduction de la
surface d'attaque obtenue.
